% section6_belief_loop_fixed_point.tex

\section{Belief Loops and Fixed-Point Selection}
\label{sec:belief-loop}

Sections~\ref{sec:surprise}--\ref{sec:dynamic} isolate the two local forces behind the pooling reversal. First, adverse repricing is pinned to true type: \(\Phi_H>\Phi_L\). Second, when the transition probability into high scrutiny is sufficiently large, D1 no longer assigns an off-path ambitious signal to the high type. This section turns those local comparisons into a fixed-point selection result.

The point is not merely that separation fails and pooling survives. That would be a linear argument. The equilibrium object is circular. Whether \(s_A\) carries a high-type premium depends on whether high types are expected to occupy \(s_A\). Whether high types occupy \(s_A\) depends on the premium they expect \(s_A\) to carry. The theorem below closes this loop.

\subsection{Belief hierarchy and psychological assessments}
\label{subsec:belief-hierarchy}

The fixed point is taken over a psychological assessment, not over a strategy profile alone. At a low-scrutiny history \(h_t\), let
\[
\sigma_\theta(h_t)\in\{s_A,s_P\}
\]
denote the type-\(\theta\) sender's signal choice. Let
\[
\rho_A(h_t):=\Pr(\theta_H\mid s_A,h_t)
\]
denote the receiver's first-order posterior after observing the ambitious signal. If \(s_A\) is on path, \(\rho_A\) is pinned down by Bayes' rule. If \(s_A\) is off path, \(\rho_A\) is selected by the refinement.

The sender's payoff depends on what she expects the receiver to believe after \(s_A\). Write
\[
\widehat \rho_{\theta,A}(h_t)
\]
for the type-\(\theta\) sender's second-order belief about the receiver's posterior following \(s_A\). A psychological assessment is a triple
\[
\mathcal A=(\sigma,\rho_A,\widehat \rho_A),
\]
where \(\widehat \rho_A=(\widehat\rho_{H,A},\widehat\rho_{L,A})\). Belief consistency requires
\[
\widehat\rho_{\theta,A}=\rho_A
\qquad\text{for each }\theta\in\{\theta_H,\theta_L\}.
\]
Thus the sender's second-order belief about the receiver's first-order posterior must agree with the posterior induced by the assessment itself.

This is the level at which the fixed point is closed. The exposure term \(\Phi_\theta\) remains true-type indexed. It is not changed by the receiver's posterior. The receiver's posterior changes the current reputational benefit of the signal; it does not relabel the deviator's adverse-repricing exposure.

Let
\[
B_\theta(\rho)
\]
denote the type-\(\theta\) current net benefit from choosing \(s_A\) rather than \(s_P\) when the ambitious signal is interpreted with posterior \(\rho\). This term includes the current reputational return from the receiver's interpretation, net of current signal cost and low-scrutiny epistemic cost. It excludes future high-scrutiny exposure. The fixed \(B_H\) used in Section~\ref{sec:dynamic} is the full high-type-attribution value:
\[
B_H=B_H(1).
\]
Assume
\[
B_H(\rho)\text{ is weakly increasing in }\rho,
\qquad
B_H(1)>0,
\qquad
B_H(0)\le 0.
\]
The first inequality says that a higher high-type posterior makes the ambitious signal more valuable. The last inequality says that an ambitious signal interpreted as low-type evidence carries no high-type premium.

For the low type, assume
\[
B_L(0)<0.
\]
Thus if \(s_A\) is interpreted as low-type evidence, the low type does not receive enough current benefit to justify choosing it.

\subsection{Depreciating exposure and the withdrawal threshold}
\label{subsec:depreciating-exposure}

The dynamic withdrawal motive requires exposure to be persistent but not permanently locked. A sender must be able to reduce future exposure by withdrawing from the ambitious signal. To make this explicit, let
\[
e_t\ge 0
\]
denote the sender's exposed stock at the beginning of period \(t\). Sending \(s_A\) adds one unit of new exposed ambition; choosing \(s_P\) adds no new unit. Existing exposure depreciates at rate \(\eta \in (0,1]\):
\[
e_{t+1}=(1-\eta)e_t+\mathbf{1}\{s_t=s_A\}.
\]
The binary exposure specification in Section~\ref{sec:model} is the special case \(\eta=1\), where withdrawal immediately removes the reduced-form exposure state. The general stock formulation makes clear why withdrawal has dynamic value: it stops new exposure from being created and lets old exposure decay.

Let
\[
\Gamma(q,\eta)>0
\]
denote the discounted marginal continuation price of adding one unit of exposed ambition in the current low-scrutiny state. In the one-period reduced form used in Section~\ref{sec:dynamic},
\[
\Gamma(q,\eta)=\beta q\mu_H.
\]
In the stock version, \(\Gamma(q,\eta)\) is the expected discounted value of the future judgment intensities applied to the marginal unit of exposure:
\[
\Gamma(q,\eta)
=
\mathbb E_q\left[
\sum_{\tau\ge 1}
\beta^\tau (1-\eta)^{\tau-1}\mu_{j_{t+\tau}}
\,\middle|\,j_t=L
\right].
\]
The maintained monotonicity condition is
\[
\Gamma(q,\eta)\text{ is continuous and strictly increasing in }q.
\]
This condition says that a higher probability of entering high scrutiny raises the continuation price of current exposure.

Given posterior \(\rho\), the high type's payoff gain from choosing \(s_A\) rather than \(s_P\) is
\[
\Delta_H(\rho,q)
=
B_H(\rho)-\Gamma(q,\eta)\Phi_H.
\]
Define the high-type withdrawal threshold by
\[
q_{\mathrm{IC}}
=
\inf\{q:\,B_H(1)-\Gamma(q,\eta)\Phi_H<0\}.
\]
In the one-period reduced form, this collapses to the threshold in Section~\ref{sec:dynamic}:
\[
q_{\mathrm{IC}}
=
\frac{B_H(1)}{\beta\mu_H\Phi_H}.
\]

\begin{lemma}[Monotone stopping]
\label{lem:monotone-stopping}
Suppose \(B_H(\rho)\) is weakly increasing in \(\rho\), \(\Phi_H>0\), and \(\Gamma(q,\eta)\) is continuous and strictly increasing in \(q\). Then the high type's stay advantage \(\Delta_H(\rho,q)\) is strictly decreasing in \(q\). In particular, if \(q>q_{\mathrm{IC}}\), then
\[
\Delta_H(\rho,q)<0
\qquad\text{for every }\rho\in[0,1].
\]
Thus the high type's withdrawal region is an upper interval in \(q\).
\end{lemma}

\begin{proof}
For fixed \(\rho\),
\[
\Delta_H(\rho,q)=B_H(\rho)-\Gamma(q,\eta)\Phi_H.
\]
Since \(\Phi_H>0\) and \(\Gamma(q,\eta)\) is strictly increasing in \(q\), \(\Delta_H(\rho,q)\) is strictly decreasing in \(q\). Since \(B_H(\rho)\le B_H(1)\) for every \(\rho\in[0,1]\), the inequality
\[
B_H(1)-\Gamma(q,\eta)\Phi_H<0
\]
implies
\[
B_H(\rho)-\Gamma(q,\eta)\Phi_H<0
\]
for every \(\rho\in[0,1]\). Hence, once \(q>q_{\mathrm{IC}}\), no receiver posterior can make the high type willing to stay on \(s_A\).
\end{proof}

Lemma~\ref{lem:monotone-stopping} is the threshold structure needed for the dynamic interpretation. The stopping region is not fragmented because the only continuation object that changes with \(q\) is the marginal exposure price, and that price moves monotonically.

\subsection{D1 and the off-path belief map}
\label{subsec:off-path-belief-map}

The next object is the receiver's belief after an off-path ambitious signal. In a pooling candidate on \(s_P\), the signal \(s_A\) is not observed on path. The receiver must decide whether an off-path \(s_A\) is more plausibly sent by \(\theta_H\) or by \(\theta_L\).

Let the off-path cost of sending \(s_A\) for type \(\theta\) be
\[
K_\theta(q)=\bar k_\theta+\Gamma(q,\eta)\Phi_\theta,
\]
where \(\bar k_\theta\) is the current off-path cost component and \(\Gamma(q,\eta)\Phi_\theta\) is the continuation exposure cost. Given receiver response \(a\), let \(R(a)\) be the reputational return from that response. The deviation gain is
\[
G_\theta(a;q)=R(a)-K_\theta(q),
\]
and the profitable-deviation response set is
\[
\mathcal D_\theta(q)=\{a:G_\theta(a;q)>0\}.
\]

Because \(\Phi_H>\Phi_L\), the high type's off-path cost rises faster with \(q\). Define
\[
q_{\mathrm{D1}}
=
\inf\{q:\,K_H(q)>K_L(q)\}.
\]
For \(q>q_{\mathrm{D1}}\),
\[
K_H(q)>K_L(q),
\]
and therefore
\[
\mathcal D_H(q)\subsetneq \mathcal D_L(q)
\]
under the strict-response richness condition used in Lemma~\ref{lem:d1-flip}. D1 then assigns the off-path ambitious signal to the low type. Hence the D1-selected posterior is
\[
\rho_A=0.
\]

The on-path and off-path belief map can therefore be summarized as follows. Let
\[
z\in\{0,1\}
\]
denote whether the high type is expected to occupy the ambitious signal:
\[
z=1 \quad\text{means } \theta_H \text{ chooses }s_A,
\]
while
\[
z=0 \quad\text{means } \theta_H \text{ withdraws to }s_P.
\]
For \(q>q_{\mathrm{D1}}\), the receiver's belief about \(s_A\) is
\[
\rho_A(z)=
\begin{cases}
1, & \text{if }z=1\text{ and }s_A\text{ is the separating high-type signal},\\
0, & \text{if }z=0\text{ and }s_A\text{ is off path, selected by D1}.
\end{cases}
\]
The first line is Bayes' rule on the separating path. The second line is D1 off path.

\subsection{The belief-loop operator}
\label{subsec:belief-loop-operator}

Given the belief map \(\rho_A(z)\), the high type's participation response is
\[
T_q(z)
=
\mathbf{1}\left\{
B_H(\rho_A(z))-\Gamma(q,\eta)\Phi_H\ge 0
\right\}.
\]
A high-type belief-loop fixed point is a value \(z\in\{0,1\}\) such that
\[
z=T_q(z).
\]
This is the self-referential equation. High-type participation determines whether \(s_A\) is on path. The path status of \(s_A\) determines the receiver's posterior. The receiver's posterior determines the current rent of \(s_A\). That rent determines high-type participation.

Define
\[
q^*=\max\{q_{\mathrm{IC}},q_{\mathrm{D1}}\}.
\]

\begin{theorem}[Belief-loop fixed-point selection]
\label{thm:belief-loop-fixed-point}
Suppose:
\begin{enumerate}
\item \(B_H(\rho)\) is weakly increasing in \(\rho\), with \(B_H(1)>0\) and \(B_H(0)\le 0\);
\item \(B_L(0)<0\);
\item \(\Phi_H>\Phi_L\), as derived from primitive self-accuracy levels and Bayesian-surprise repricing;
\item \(\Gamma(q,\eta)\) is continuous and strictly increasing in \(q\), with \(\eta \in (0,1]\);
\item \(q>q^*=\max\{q_{\mathrm{IC}},q_{\mathrm{D1}}\}\);
\item off-path beliefs satisfy D1;
\item the analysis is restricted to the binary-signal pure-strategy class.
\end{enumerate}
Then the unique D1-consistent psychological assessment in this class has
\[
\theta_H\mapsto s_P,
\qquad
\theta_L\mapsto s_P.
\]
Equivalently, the unique high-type belief-loop fixed point is
\[
z^*=0.
\]
\end{theorem}

\begin{proof}
First consider the candidate in which the high type occupies the ambitious signal, \(z=1\). In the separating interpretation, \(s_A\) is on path for the high type, so Bayes' rule gives
\[
\rho_A(1)=1.
\]
The high type's payoff gain from choosing \(s_A\) rather than \(s_P\) is
\[
\Delta_H(1,q)=B_H(1)-\Gamma(q,\eta)\Phi_H.
\]
Since \(q>q_{\mathrm{IC}}\), this expression is strictly negative. Hence
\[
T_q(1)=0.
\]
Therefore \(z=1\) is not a fixed point.

Now consider the candidate in which the high type withdraws, \(z=0\). Then \(s_A\) is off path for the high type. Since \(q>q_{\mathrm{D1}}\), the D1 flip assigns an off-path \(s_A\) to the low type. Hence
\[
\rho_A(0)=0.
\]
The high type's payoff gain from deviating to \(s_A\) is
\[
\Delta_H(0,q)=B_H(0)-\Gamma(q,\eta)\Phi_H.
\]
Since \(B_H(0)\le 0\) and \(\Gamma(q,\eta)\Phi_H>0\),
\[
\Delta_H(0,q)<0.
\]
Thus
\[
T_q(0)=0.
\]
Therefore \(z=0\) is a fixed point.

It remains to rule out pure-strategy profiles in which the low type occupies \(s_A\). If only the low type chooses \(s_A\), then \(s_A\) is on path as low-type evidence, so \(\rho_A=0\). The low type's gain is
\[
B_L(0)-\Gamma(q,\eta)\Phi_L<0,
\]
because \(B_L(0)<0\) and \(\Gamma(q,\eta)\Phi_L>0\). Hence the low type does not wish to occupy \(s_A\) under low-type attribution.

If both types choose \(s_A\), then \(s_A\) is pooling rather than separating. Let \(\pi_0\in(0,1)\) be the prior probability of \(\theta_H\). Bayes' rule gives \(\rho_A=\pi_0\). Since \(B_H(\rho)\) is weakly increasing and \(\pi_0<1\),
\[
B_H(\pi_0)\le B_H(1).
\]
Because \(q>q_{\mathrm{IC}}\),
\[
B_H(1)-\Gamma(q,\eta)\Phi_H<0.
\]
Therefore
\[
B_H(\pi_0)-\Gamma(q,\eta)\Phi_H<0.
\]
The high type would deviate from pooling on \(s_A\) to the safe signal. Hence pooling on \(s_A\) is not a pure-strategy equilibrium.

The only remaining pure-strategy profile is pooling on \(s_P\). Under that profile, \(s_A\) is off path. D1 gives \(\rho_A=0\), and the two deviation inequalities above show that neither type gains from deviating to \(s_A\). Thus pooling on \(s_P\) is the unique D1-consistent psychological assessment in the binary pure-strategy class.
\end{proof}

\subsection{The closed loop}
\label{subsec:closed-loop}

The theorem closes the loop in both directions. If high types are expected to stay, then \(s_A\) is read as high type:
\[
z=1\Rightarrow \rho_A=1.
\]
But when \(q>q_{\mathrm{IC}}\), even full high-type attribution is not enough to offset continuation exposure:
\[
\rho_A=1\Rightarrow T_q(1)=0.
\]
Thus the separating belief system destroys itself:
\[
z=1\Rightarrow \rho_A=1\Rightarrow T_q(1)=0\ne z.
\]

If high types are expected to withdraw, then \(s_A\) is off path. When \(q>q_{\mathrm{D1}}\), D1 assigns the off-path ambitious deviation to the low type:
\[
z=0\Rightarrow s_A\text{ off path}\Rightarrow \rho_A=0.
\]
Once \(\rho_A=0\), the ambitious signal carries no high-type premium, while it still creates true-type exposure for a high-type deviator:
\[
\rho_A=0\Rightarrow T_q(0)=0.
\]
Thus the withdrawal belief reproduces itself:
\[
z=0\Rightarrow \rho_A=0\Rightarrow T_q(0)=0=z.
\]

This is the belief loop. High types withdraw because exposed ambition is too costly. Their withdrawal makes \(s_A\) off path. D1 then removes the high-type attribution from \(s_A\). Once that attribution is removed, the current rent from returning to \(s_A\) disappears. Withdrawal is therefore self-confirming.

\subsection{No interior mixing in this section}
\label{subsec:no-mixing-sec6}

The fixed-point theorem is a pure-strategy selection result. It does not use an interior mixing probability for the high type. This restriction is deliberate. In a binary model with no low-type background arrival at \(s_A\), any strictly positive high-type mass on \(s_A\) would keep the posterior after \(s_A\) degenerate in the high type. Continuous signal dilution therefore cannot be generated by high-type mixing alone.

Interior mixing requires an additional object, such as low-type mimetic entry, noisy audience inference, trembles, or a continuous type distribution. Those mechanisms belong to the mixed-signal extension. The present section proves only the sharper fixed-point claim: above \(q^*\), the only D1-consistent pure belief loop is full withdrawal from the ambitious signal and pooling on \(s_P\).
